\subsection{Does a learned refiner help?}
\label{sec:exp_refiner}
% Learned refiner on the 526 clips, all rows read in one evaluation job (workbook sheet 31, complete version).
% W/L vs the same front end without the refiner. Appearance normalised within this job.
\begin{table}[t]
\centering
\caption{\textbf{Does a learned refiner help?} The 1.3B bridge-flow refiner applied on top of \method's front end (paste only) and of the box policy, with and without the carrier-anchored attention bias ($\lambda=2$). 526 clips, one evaluation job; W/L against the same front end without the refiner. Boundary jump: change at the refiner's 17-frame chunk seams relative to other frame pairs (1 = no seam).}
\label{tab:refiner}
\resizebox{\linewidth}{!}{%
\begin{tabular}{@{}lcccccc@{}}
\toprule
 & Read.\ $\uparrow$ & W/L & Flicker $\downarrow$ & Jitter $\downarrow$ & B.\ jump $\downarrow$ & Appear.\ $\downarrow$ \\
\midrule
\steer & 84.2 & -- & 1.66 & 8.34 & 0.93 & 1.58 \\
\midrule
\method front end (paste only) & 91.3 & -- & 1.77 & 7.69 & 0.84 & 2.49 \\
\quad + refiner & 87.8 & 3/21 & 1.96 & 8.22 & 2.26 & 2.45 \\
\quad + refiner + attention bias & 88.6 & 3/17 & 2.00 & 8.41 & 2.06 & 2.33 \\
\midrule
Box policy (paste only) & 91.4 & -- & 1.77 & 7.83 & 0.84 & 2.35 \\
\quad + refiner & 87.6 & 2/22 & 1.99 & 8.46 & 2.31 & 2.31 \\
\quad + refiner + attention bias & 88.4 & 3/19 & 2.03 & 8.68 & 2.10 & 2.24 \\
\bottomrule
\end{tabular}}
\end{table}


Rendered glyphs look pasted, so the obvious remedy is to let a generative model harmonise them with the scene while keeping their shapes. We invested heavily in this route before settling on the training-free pipeline, and report it in full because its outcome shapes our design.

\paragraph{Refiner.} The refiner is a video diffusion transformer initialised from Wan2.1-T2V-1.3B~\cite{wan2025wan} that takes the pasted clip as extra input channels (zero-initialised) and is trained as a flow-matching bridge~\cite{lipman2023flowmatching,liu2023rectifiedflow} from the paste to the target; inference starts from the paste, in the spirit of SDEdit~\cite{meng2022sdedit}, and takes 30 Euler steps on 17-frame chunks of $384\times384$ crops. A glyph loss compares Sobel edges of the decoded prediction with the target inside a band around the glyph (weighted by the signed distance to the glyph boundary). Training pairs come from three sources: a Blender data factory with six carrier surfaces and passing occluders, which renders every sample twice (lit composite and flat paste); 562 real pairs obtained by erasing correct text from generated clips and re-pasting it; and up to 1.15 million image pairs built the same way from AnyWord-3M (63\% LAION, 31\% Wukong, 6\% street scenes). All refiners were trained on H100-class GPUs (GH200/H200) for 20k steps.

\paragraph{More data is not the lever.} On a 296-clip garbled subset, refiners trained with 0, 10k, 100k and 1M image pairs reach 78.4, 75.3, 76.4 and 76.4\% readability, all below the paste they start from (80.7\%), and their appearance does not improve with data (supplement). A diagnostic on 1{,}500 held-out pairs explains why: on the very task they were trained for---turning the paste back into the original---all of them move \emph{away} from the original, they do not even memorise pairs seen twelve times, and they fail to learn a systematic difference (originals have thicker strokes than our pastes in 60\% of the pairs). The supervision was inconsistent: the paste uses our font, the target uses the generator's, and the glyph loss, which dominates the objective, effectively teaches the model to erase our strokes. Making the supervision glyph-consistent---pasting the original's own ink mask instead of our font---is the first recipe that moves toward the original on held-out pairs and it cuts the image-domain readability loss from 26 to 4 points, with appearance better than the paste. On videos, however, it still does not exceed the paste (best variant: 79.7 vs.\ 80.7\%).

\paragraph{System-level effect.} \Cref{tab:refiner} applies the best refiner to the full benchmark. On top of \method's front end it costs 3.5 points of readability (21 clips broken against 3 repaired; $p<0.001$), increases flicker and jitter, and introduces seams at its chunk boundaries (boundary jump 2.26 vs.\ 0.84), while the appearance distance does not change significantly ($-0.035$, 95\% CI $[-0.21, 0.19]$). A training-free attention bias that ties tokens to their carrier trajectory, in the spirit of flow-guided attention~\cite{cong2024flatten,geyer2024tokenflow}, was selected on 50 clips with an earlier refiner (boundary jump $-16\%$ there) and reduces the seams by 9\% on the full set (2.06 vs.\ 2.26; better on 364 of 526 clips), but it slightly increases flicker within chunks and leaves the readability loss in place. We therefore keep the refiner out of \method. The negative result is specific to our scale---1.3B parameters and at most a few hundred real video pairs, against a 14B backbone and a million training triplets for \steer---but it is consistent: at this scale, a learned harmoniser trades the guarantee that motivates rendering for an appearance gain we could not measure.
